% =====================================================================
% Grad-CAM, lung masks, attention metric and shortcut findings  (author: Aman Singh, Stage 3)
% Every number below comes from results/ in the repository:
%   results/stage3_metrics.csv, results/stage3/attention_summary.csv, results/stage3/attention_robustness.csv,
%   results/stage3/lung_mask_summary.csv, results/stage3/sample_images.csv, results/stage3/settings.json
% Figures to upload to Overleaf (keep these paths):
%   figures/stage3/lung_masks_random.png          <- results/stage3/lung_masks_random.png
%   figures/stage3/attention_boxplot.png          <- results/stage3/attention_boxplot.png
%   figures/stage3/false_positives_densenet121.png <- results/stage3/false_positives_densenet121.png
% Labels used from other sections: sec:data-preproc, sec:data-eda, sec:densenet, tab:comparison
% Citation keys: selvaraju2017gradcam, geirhos2020shortcut, zech2018variable, degrave2021ai (references_abhishek.bib);
%   zhao2017pspnet, lian2021chestxdet, wang2017chestxray8, cohen2022torchxrayvision, wang2020scorecam (references_aman.bib)
% =====================================================================
\section{Where Do the Models Look? Grad-CAM and Lung Attention}
\label{sec:gradcam}

Both classifiers reach a high test AUC (\cref{tab:comparison}), but a high score does not show that a model uses the lungs.
In this section we measure where the custom CNN and DenseNet-121 look, and how much of that attention falls on the lungs.

\subsection{Grad-CAM for a single-logit classifier}
\label{sec:gradcam-method}
Grad-CAM~\cite{selvaraju2017gradcam} explains a decision with the feature maps $A^k$ ($k=1,\dots,K$) of a late convolutional layer.
Each map is weighted by the average gradient of the explained score $s$, and the weighted sum is passed through a ReLU:
\begin{equation}
  \alpha_k = \frac{1}{hw}\sum_{i,j}\frac{\partial s}{\partial A^k_{ij}}, \qquad
  L = \mathrm{ReLU}\Big(\sum_k \alpha_k A^k\Big).
  \label{eq:gradcam}
\end{equation}
Our models have one output logit $z$ with $\sigma(z)=P(\text{pneumonia})$ and no separate NORMAL output.
We therefore explain the \emph{predicted} class: $s=z$ when the image is predicted PNEUMONIA ($\sigma(z)\ge0.5$) and $s=-z$ when it is predicted NORMAL,
so every heatmap shows the evidence for the class the model chose.
We hook the activations after the last ReLU, which are the features the classifier actually receives:
\texttt{features.block5.relu2} for the custom CNN ($256\times14\times14$) and the post-ReLU feature map of DenseNet-121 ($1024\times7\times7$).
The map $L$ is upsampled bilinearly to $224\times224$, the input size, so it lines up pixel by pixel with the lung masks.
A toy model with a known answer confirms that the heatmap lands on the right image region for both predicted classes.

\subsection{Lung masks}
\label{sec:gradcam-masks}
We segment the lungs of all 5,856 images with the pretrained PSPNet~\cite{zhao2017pspnet} from TorchXRayVision~\cite{cohen2022torchxrayvision},
which was trained on the ChestX-Det annotations~\cite{lian2021chestxdet} of adult NIH chest X-rays~\cite{wang2017chestxray8}.
We join its left- and right-lung outputs into one mask.
The classifiers see the whole image squashed to $224\times224$ (\cref{sec:data-preproc}), so we also resize the whole image without cropping before segmentation,
and resize the mask back to $224\times224$ with nearest-neighbour interpolation to keep it binary.
\Cref{fig:masks} shows masks for 24 random images.
% TODO (Aman): after the eye check of results/stage3/lung_masks_random.png, state how many of the 24 masks look right
% (the failures seen so far are pneumonia images with a partly missed, consolidated lung).

The segmentation model was not trained on children, and it fails on some pneumonia images: a lung that is white from consolidation is partly or entirely missed.
We flag a mask when the lungs cover less than 8\% or more than 60\% of the image, or when one lung is under 5\% of the image or under 40\% of the other lung.
Of the 5,856 masks, 226 are flagged; 225 of them are PNEUMONIA images (181 train, 25 validation, 19 test).

\begin{figure*}[t]
  \centering
  \includegraphics[width=\linewidth]{figures/stage3/lung_masks_random.png}
  \caption{Lung masks (red, outlined in green) for 24 random images, drawn on the $224\times224$ input the classifiers see. Most masks cover both lung fields; failures occur on pneumonia images where consolidated lung is missed.}
  \label{fig:masks}
\end{figure*}

\subsection{Lung-attention metric}
\label{sec:gradcam-metric}
For a heatmap $L\ge0$ and a lung mask $M\in\{0,1\}$ of the same size, the lung attention is the share of the heatmap inside the lungs:
\begin{equation}
  \text{attention\_in\_lung} = \frac{\sum_{i,j} L_{ij} M_{ij}}{\sum_{i,j} L_{ij}}.
  \label{eq:attention}
\end{equation}
This is the energy-based pointing game of Score-CAM~\cite{wang2020scorecam} with the lung mask in place of a bounding box.
The value alone is hard to read: a heatmap spread evenly over the image would score the lung area $\frac{1}{hw}\sum_{i,j}M_{ij}$, about 38\% on our test images.
We therefore report the lung area as the chance level next to each result, and the \emph{lift}, the ratio of the two, where a lift above 1 means more attention on the lungs than chance.
We also report the \emph{border share}, the part of the heatmap in the outer 10\% of the image on each side, where text markers and scanner labels sit;
an even heatmap has a border share of $1-0.8^2=36\%$.
All values are computed on the 624 test images for both models.

\subsection{Results}
\label{sec:gradcam-results}
\begin{table}[t]
  \centering
  \caption{Lung attention on the test set (mean over images). Lung area is the chance level: the score of an evenly spread heatmap. Lift $>1$ is the share of images with more attention on the lungs than chance.}
  \label{tab:attention}
  \small
  \begin{tabular}{llrrrr}
    \toprule
    Model & Images & Attention & Lung area & Lift $>1$ & Border \\
    \midrule
    Custom CNN   & all (624)       & 43.4\% & 37.7\% & 76.8\% & 26.8\% \\
                 & NORMAL (234)    & 42.3\% & 39.4\% & 64.5\% & 27.7\% \\
                 & PNEUMONIA (390) & 44.1\% & 36.7\% & 84.1\% & 26.2\% \\
    \midrule
    DenseNet-121 & all (624)       & 39.5\% & 37.7\% & 60.1\% & 30.9\% \\
                 & NORMAL (234)    & 42.9\% & 39.4\% & 68.8\% & 28.5\% \\
                 & PNEUMONIA (390) & 37.5\% & 36.7\% & 54.9\% & 32.4\% \\
    \bottomrule
  \end{tabular}
\end{table}

\Cref{tab:attention} shows that both models put only slightly more of their heatmap on the lungs than an evenly spread heatmap would:
43.4\% for the custom CNN and 39.5\% for DenseNet-121, against a lung area of 37.7\%.
The more accurate model is the one closer to chance.
On PNEUMONIA images, DenseNet-121 puts 37.5\% of its attention on the lungs, which cover 36.7\% of the image, and only 54.9\% of these images have a lift above~1.
High accuracy therefore does not mean that the model bases its decision on the lungs.
\Cref{fig:attention-box} shows the spread per class and for correct and wrong predictions.

\begin{figure*}[t]
  \centering
  \includegraphics[width=\linewidth]{figures/stage3/attention_boxplot.png}
  \caption{Lung attention per test image, by true class and correct or wrong prediction. The red dash is the mean lung area of each group, the chance level.}
  \label{fig:attention-box}
\end{figure*}

\subsection{Shortcut examples}
\label{sec:gradcam-shortcuts}
The clearest shortcut appears in DenseNet-121's most confident mistakes (\cref{fig:shortcuts}).
For NORMAL test images that it predicts as PNEUMONIA with probability 1.00, the heatmap concentrates on text burned into the image:
the \texttt{KV58}/\texttt{KV59}/\texttt{KV60} exposure labels in the bottom-right corner and the timestamps in the top-right corner, not the lungs.
The numbers agree with the pictures. On NORMAL images that are wrongly predicted as PNEUMONIA, the border share is 35.8\% for the custom CNN (19.6\% on correctly predicted NORMAL images)
and 34.6\% for DenseNet-121 (25.4\% on correct ones).
The custom CNN's heatmaps often sit on the diaphragm and the strip below the lungs.
These findings match the shortcuts reported for chest X-ray classifiers by Zech~\etal~\cite{zech2018variable} and DeGrave~\etal~\cite{degrave2021ai},
and the general pattern of shortcut learning described by Geirhos~\etal~\cite{geirhos2020shortcut}. They also add to the recording differences between the classes found in the EDA (\cref{sec:data-eda}).
% TODO (Aman): add the 6-8 hand-picked examples of Stage 3.4 (issue #14) from results/stage3/ and describe each one.

\begin{figure*}[t]
  \centering
  \includegraphics[width=\linewidth]{figures/stage3/false_positives_densenet121.png}
  \caption{DenseNet-121's most confident false positives (NORMAL test images predicted as PNEUMONIA). The heatmaps focus on burned-in text such as the \texttt{KV58}/\texttt{KV59}/\texttt{KV60} exposure labels (bottom right) and timestamps (top right) instead of the lungs. Green: lung mask outline.}
  \label{fig:shortcuts}
\end{figure*}

\subsection{Limitations}
\label{sec:gradcam-limits}
The lung masks come from a model trained on adult X-rays and can miss consolidated lung, so attention on diseased lung may be counted as outside the lungs.
This affects PNEUMONIA images only. Leaving out the 19 flagged test images changes the mean lung attention by less than 0.01 for every model and class (largest change: custom CNN on PNEUMONIA images, 44.1\% to 44.7\%),
but masks that are only partly wrong are not all flagged.
DenseNet-121's heatmaps are computed on a $7\times7$ map (the custom CNN's on $14\times14$), and coarser maps spread out more after upsampling, which moves the lung attention towards the chance level.
Grad-CAM shows which regions raised the score, not why, and it does not prove that the model would fail without them;
\cref{sec:masking} tests this directly by training on lung-masked images. % TODO (Himanshu): check the label of the Stage 4 section
